% id: 109-20
% book: 베이직쎈 확률과 통계 · 06 이항분포와 정규분포
% section: 개념 32 정규분포곡선에서의 확률
% passage: 확률변수 $X$가 정규분포 $\mathrm{N}(m{,}\mkern3mu\mathopen{}\sigma^2)$을 따르고\par\smallskip\centerline{$\mathrm{P}(m\le X\le m+\sigma)=a$, $\mathrm{P}(m\le X\le m+2\sigma)=b$}\par\smallskip\noindent 일 때, 다음을 $a$, $b$를 사용하여 나타내시오.
% figure: none
% answer: 
% answer_source: 
% variant_level: 0 (원본 전사)
% 이 파일은 build.mjs 가 items.json 에서 만든다. 고칠 때는 items.json 을 고친다.
\smash{\raisebox{1.5mm}{\dmkichul{베이직쎈 확률과 통계 109쪽 20번}}}
$\mathrm{P}(m-\sigma\le X\le m+2\sigma)$
